\documentclass[../main.tex]{subfiles}
\graphicspath{{\subfix{../image/}}}

\begin{document}

\chapter{指数映射}\label{chapter:ch20-exponential-map}

本章中，我们应用流、Lie 导数与叶状结构的工具，深入探讨 Lie 群与 Lie 代数之间的关系。

第一节中，我们定义 Lie 群 $G$ 的单参数子群（即从 $\mathbb{R}$ 到 $G$ 的 Lie 群同态），并证明 $\operatorname{Lie}(G)$ 的元素与 $G$ 的单参数子群之间存在一一对应。

接着引入我们研究的焦点，即从 Lie 代数到群的一个典范光滑映射，称为指数映射。它把 $\operatorname{Lie}(G)$ 中过原点的直线映为 $G$ 的单参数子群。

作为指数映射的第一个主要应用，我们证明闭子群定理，它说 Lie 群的每个拓扑闭子群实际上都是嵌入 Lie 子群。

随后我们证明流基本定理的一个高维推广。不再是单个光滑向量场生成 $\mathbb{R}$ 的作用，而是考虑一族有限维向量场，追问它们何时生成某个 Lie 群的作用。主要定理是：若 $G$ 是单连通 Lie 群，则任意从 $\operatorname{Lie}(G)$ 到 $M$ 上完全向量场集合的 Lie 代数同态都生成 $G$ 在 $M$ 上的一个光滑作用。

最后两节中，我们把所有这些结果汇集起来，加深对 Lie 群与 Lie 代数之间对应的理解。首先证明有限维 Lie 代数的同构类与单连通 Lie 群的同构类之间存在一一对应；然后证明对任意 Lie 群 $G$，$G$ 的连通法向子群对应于 $G$ 的 Lie 代数中的理想，即在与代数中任意元素作括号运算下保持稳定的子空间。这极好地诠释了 Lie 理论的基本哲学：尽可能利用 Lie 群/Lie 代数对应，把关于 Lie 群的群论问题翻译为关于其 Lie 代数的线性代数问题。

\subfile{Chapter-20/section-01.tex}
\subfile{Chapter-20/section-02.tex}
\subfile{Chapter-20/section-03.tex}
\subfile{Chapter-20/section-04.tex}
\subfile{Chapter-20/section-05.tex}
\subfile{Chapter-20/section-06.tex}

\end{document}
